\silentchapter{Proposisjonar}{proposisjonar}

\prop{1}{}{Formverda er alt som er tilfelle.}
\prop{1.1}{o}{Formverda er totaliteten av realiserte konfigurasjonar, ikkje ting\footnote{Ontologien likestiller fysiske samanføyningar, distribuerte nettverk og sekvensielle systemkall.}.}
\prop{1.11}{t}{Alt som har form, har eksakt denne forma\footnote{Maskintilstanden fastslår identiteten sin absolutt i eksekveringsaugnen.}.}
\prop{1.12}{t}{At ein konfigurasjon tek éi form framfor andre moglege, krev ei ekstern årsak\footnote{Minneallokeringa peikar utover sine eigne byte, mot prosessen som kravde ho.}.}
\prop{1.121}{t}{Årsaka kviler i relasjonen mellom den manifesterte konfigurasjonen og dei urealiserte moglegheitene\footnote{Eit grensesnitt får tyngde ved å utelukke latente interaksjonsstiar.}.}
\prop{1.1211}{t}{Dette er ein relasjonell, ikkje absolutt, eigenskap\footnote{Signalet får informasjonsverdi utelukkande frå rommet av usende signal.}.}
\prop{1.2}{d}{Eit \textbf{objekt}\footnote{Databasens atomære oppføring; asynkronarkitekturens udelelege tilstandspakke.} er den enklaste bestanddelen i ei form.}
\prop{1.21}{t}{Objektet er udeleleg og uforanderleg\footnote{Dekomponering under dette nivået oppløyser funksjonen og sprengjer handlingsrommets ontologi.}.}
\prop{1.211}{t}{Objekta utgjer substansen i formverda\footnote{Primitivrepertoaret definerer yttergrensa for kva systemet kan kompilere.}.}
\prop{1.2111}{t}{Alle moglege former er alt gjevne med objekta\footnote{Kryptoprotokollen skisserer transaksjonstopologien før første køyring.}.}
\prop{1.21111}{o}{Moglegheita eksisterer logisk fordi objekta finst, uavhengig av tenking\footnote{Parseren evaluerer syntakstreet blindt for programmerarens intensjon.}.}
\prop{1.22}{d}{Ein \textbf{konfigurasjon}\footnote{Nodegrafen bestemmer systemeigenskapane strengt, uavhengig av domene.} er ei spesifikk samansetjing av objekt.}
\prop{1.221}{t}{Forma er strukturen til konfigurasjonen\footnote{Substratet kan skifte frå silisium til karbon; overføringsfunksjonen består.}.}
\prop{1.3}{d}{\textbf{Formrommet}\footnote{Søkerommet til ein genetisk algoritme er isomorft med dei tillatne morfologiske tilstandane.} (\textit{morphospace}) til ein klasse er mengda av alle hennar moglege konfigurasjonar.\footnote{Raup (1966); Mitteroecker \& Huttegger (2009)}}
\prop{1.31}{d}{Kvart realisert objekt okkuperer eitt eksakt punkt i dette n-dimensjonale rommet\footnote{Tilstandsvektoren gjev konfigurasjonen ein irreversibel og eintydig indeks.}.}
\prop{1.32}{d}{Det empiriske formrommet er alltid ein projeksjon\footnote{Kompresjon av høgdimensjonale data til metrikkar tvingar fram oppløysingstap.}.}
\prop{1.321}{t}{Inga endeleg parametermengd fangar den latente kompleksiteten fullt ut.\footnote{Thompson (1917)}}
\prop{1.3211}{o}{Projeksjonen vel: ho lyser opp somme relasjonar og mørklegg andre\footnote{Funksjonell optimalisering degenererer samstundes unytta kapabilitetar.}.}
\prop{1.32111}{o}{Det finst ingen nøytral projeksjon\footnote{Aggregering brenn bort informasjon; restane speglar eit vilkårleg føremål.}.}
\prop{1.32112}{o}{Projeksjonen gjer empirien mogleg, men konstituerer ikkje forma\footnote{Skjermteikninga er projeksjonen; minnetreet ber den logiske forma.}.}
\prop{1.321121}{o}{Forma eksisterer uavhengig av målinga\footnote{Tilstanden eksisterer uavhengig av observatøren.}.}
\prop{1.33}{d}{Klassen avgrensar formrommet funksjonelt\footnote{Typesystemet rammar inn operasjonane; det uvaliderbare fell utanfor.}.}
\prop{1.331}{t}{Klassen set grensene; seleksjonstrykket fordelar innhaldet\footnote{Bandbreidde og merksemd filtrerer kva tilstandar systemet akkumulerer.}.}
\prop{1.4}{d}{Formrommet deler seg topologisk i tre regionar: busette, opne og forbodne\footnote{Design styrer grensesnitta mellom desse tre ontologiske fasane.}.}
\prop{1.41}{o}{Dei busette regionane utgjer det historiske arkivet\footnote{Versjonsloggen rommar heile slektstreet av levedyktige mutasjonar.}.}
\prop{1.411}{o}{Dei fungerer som ankerpunkt for framtidig navigasjon\footnote{Konvensjonar herdar til faste stasjonar systemet lyt passere.}.}
\prop{1.4111}{t}{Tyngda til ankerpunktet veks med okkupasjonstid og konvergerande tradisjonar\footnote{Standardiserte protokollar skapar gravitasjonsfelt som tvingar nye nodar i bane.}.}
\prop{1.41111}{t}{Uavhengig konvergens provar eit reelt optimum sterkare enn éin einskild tradisjon\footnote{Når urelaterte modellar finn same heuristikk, speglar dette ei lovmessigheit, ikkje ein treningsartefakt.}.}
\prop{1.42}{o}{Dei opne regionane utgjer det \textit{tilstøytande moglege}\footnote{Dette utgjer den presise innovasjonsradiusen systemet toler i neste iterasjon.}.}
\prop{1.421}{o}{Dei er tilgjengelege, men uaktualiserte\footnote{Avkutta kodespor kviler kompilert i basen, uoppkalla før eit ytre mønster bryt straumen.}.}
\prop{1.4211}{t}{Det tilstøytande moglege er endeleg\footnote{Det kombinatoriske rommet krympar asymmetrisk når systemintegrasjonen aukar.}.}
\prop{1.42111}{t}{Det er spent ut mellom busette naboar og forbodne barrierar\footnote{Migrering støyter asymmetrisk mot eldre, hardkoda avhengigheiter.}.}
\prop{1.42112}{o}{Det busette tiltrekkjer; det forbodne fråstøyter\footnote{Brukarvennlegheit tiltrekkjer; kognitiv overbelasting utgjer ein frastøytande horisont.}.}
\prop{1.421121}{o}{Det realiserte rommet oppstår i spennet mellom dei to\footnote{Arkitekturen stabiliserer seg i likevekta mellom latens og gjennomstrøyming.}.}
\prop{1.43}{d}{Dei forbodne regionane er utelukka av materielle, strukturelle eller energetiske restriksjonar\footnote{Konkurransen i transaksjonsminnet definerer uframkallelege asymmetriar.}.}
\prop{1.431}{t}{Forbodet følgjer vilkåra, ikkje regionen\footnote{Asynkrone trådar opnar augneblikkeleg den synkront uoppnåelege arkitekturen.}.}
\prop{1.4311}{o}{Grensa mellom det opne og det forbodne er dynamisk\footnote{Kompilatoroppdateringa flyttar randsona til heile handlingsrommet.}.}
\prop{1.43111}{o}{Eit materialskifte kan opne ein forboden region\footnote{Eit lagringsparadigmeskifte inverterer kostnaden av operasjonar.}.}
\prop{1.43112}{o}{Grensa tilhøyrer vilkåra, ikkje forma\footnote{Nettverksrestriksjonane ligg i protokoll og maskinvare, ikkje i datasettet.}.}